\subsection{General Background and Summary}
{\small\begin{longtable}{p{.005\textwidth}p{.995\textwidth}}
\multicolumn{2}{l}{\textbf{bibkey}} \\
& \textit{Description}: The unique identifier to match the reviewed paper to a \texttt{.bib} file. \\
& \textit{Codebook question}: ID of article (this is provided in the list of papers) \\
\\
\multicolumn{2}{l}{\textbf{title}} \\
& \textit{Description}: The title of the article. \\
& \textit{Codebook question}: Title of article \\
\\
\multicolumn{2}{l}{\textbf{benchmark}} \\
& \textit{Description}: The name of the benchmark. \\
& \textit{Codebook question}: The name of the benchmark, if one exists (e.g. GSM8K) \\
\\
\multicolumn{2}{l}{\textbf{inclusion}} \\
& \textit{Description}: Whether the paper was included in the review. \\
& \textit{Codebook question}: According to the criteria, should this paper be included or excluded? \\
\\
\multicolumn{2}{l}{\textbf{exclusion\_criteria}} \\
& \textit{Description}: The criteria for excluding the paper, if any. \\
& \textit{Codebook question}: If exclude, what criteria is violated? \\
\\
\multicolumn{2}{l}{\textbf{exclusion\_criteria\_detail}} \\
& \textit{Description}: Any additional details about the exclusion criteria. \\
& \textit{Codebook question}: If exclude, why? (optional, 1 sentence) \\
\\
\multicolumn{2}{l}{\textbf{short\_summary}} \\
& \textit{Description}: A short summary of the paper. \\
& \textit{Codebook question}: Short summary of paper contribution and method. Likely to be similar to the abstract. (2-3 sentences, no need for numbers) \\
\\
\multicolumn{2}{l}{\textbf{contribution}} \\
& \textit{Description}: Any additional notes about the article contribution \\
& \textit{Codebook question}: Other useful notes on contribution details (optional, only if something stood out) \\
\\
\end{longtable}}

\subsection{\textcolor{phenomenon}{Phenomenon}}
{\small\begin{longtable}{p{.005\textwidth}p{.995\textwidth}}
\multicolumn{2}{l}{\textbf{target\_phenomenon}} \\
& \textit{Description}: The main phenomenon measured in the paper, as defined by the authors. \\
& \textit{Codebook question}: According to the authors, what capability or specific application is being measured? (a few words, e.g. knowledge, reasoning, natural language understanding) \\
\\
\multicolumn{2}{l}{\textbf{phenomenon\_short}} \\
& \textit{Description}: Whether the phenomenon is a general capability or a specific application. \\
& \textit{Codebook question}: Which category does the target phenomenon fall into? \\
\\
& \textit{Summary of values:} \\
& General Capability (A broadly useful ability, which could be relevant to multiple applications): 321 \\
& Specific Application (A single use case, where the benchmark is likely to be examples of that use case): 118 \\
& Other: 16 \\
\\
\multicolumn{2}{l}{\textbf{phenomenon\_defined}} \\
& \textit{Description}: Whether the phenomenon is defined in the paper. \\
& \textit{Codebook question}: Is the targeted phenomenon explicitly defined? \\
\\
& \textit{Summary of values:} \\
& Yes: 348 \\
& No: 99 \\
\\
\multicolumn{2}{l}{\textbf{phenomenon\_definition}} \\
& \textit{Description}: The definition of the phenomenon. \\
& \textit{Codebook question}: How is the phenomenon of interest defined? (copy paste if possible, otherwise summarise what is being said) \\
\\
\multicolumn{2}{l}{\textbf{phenomenon\_taxonomy\_root}} \\
& \textit{Description}: The root category of the phenomenon taxonomy. \\
& \textit{Codebook question}: None \\
\\
\multicolumn{2}{l}{\textbf{phenomenon\_taxonomy\_leaf}} \\
& \textit{Description}: The leaf category of the phenomenon taxonomy. \\
& \textit{Codebook question}: None \\
\\
\multicolumn{2}{l}{\textbf{phenomenon\_taxonomy\_alternate}} \\
& \textit{Description}: An alternate for the phenomenon taxonomy if highly relevant. \\
& \textit{Codebook question}: None \\
\\
\multicolumn{2}{l}{\textbf{phenomenon\_contested}} \\
& \textit{Description}: Whether the definition of the phenomenon is broadly agreed upon, or if many definitions exist for the same term. \\
& \textit{Codebook question}: Does the target phenomenon have a widely agreed-upon definition, or is this definition contested? \\
\\
& \textit{Summary of values:} \\
& Contested: 225 \\
& Widely-agreed: 203 \\
& Not defined: 27 \\
\\
\multicolumn{2}{l}{\textbf{phenomenon\_contested\_clean}} \\
& \textit{Description}: Standardised mapping of phenomenon\_contested values for statistical analysis. \\
& \textit{Codebook question}: None \\
\\
& \textit{Summary of values:} \\
& ['Contested']: 225 \\
& ['Widely-agreed']: 203 \\
& ['No definition']: 27 \\
\\
\multicolumn{2}{l}{\textbf{definition\_scope}} \\
& \textit{Description}: Whether the benchmark covers everything within the phenomenon definition or only a subset. \\
& \textit{Codebook question}: Does the benchmark claim to measure everything covered by the definition, or focus on a more specific case or subset? \\
\\
& \textit{Summary of values:} \\
& Subset: 253 \\
& Comprehensive: 196 \\
\\
\multicolumn{2}{l}{\textbf{definition\_integrity}} \\
& \textit{Description}: Whether the definition is described as containing separate sub-phenomena. \\
& \textit{Codebook question}: Do the authors describe the phenomena as a single cohesive whole, or does it consist of sub-elements? \\
\\
& \textit{Summary of values:} \\
& Composite phenomenon: 278 \\
& Single cohesive phenomenon: 166 \\
& Other: 10 \\
\\
\multicolumn{2}{l}{\textbf{definition\_integrity\_detail}} \\
& \textit{Description}: If the definition includes sub-elements, what are they? \\
& \textit{Codebook question}: If the target phenomenon consists of sub-elements, are they measured separately? \\
\\
& \textit{Summary of values:} \\
& Yes: 261 \\
& Not applicable: 144 \\
& No: 46 \\
\\
\multicolumn{2}{l}{\textbf{purpose\_extra}} \\
& \textit{Description}: Any additional notes about the conceptual details of the paper. \\
& \textit{Codebook question}: Other useful notes on conceptual details (optional, only if something stood out) \\
\\
\end{longtable}}

\subsection{\textcolor{task}{Task} and Dataset}
{\small\begin{longtable}{p{.005\textwidth}p{.995\textwidth}}
\multicolumn{2}{l}{\textbf{task\_definition}} \\
& \textit{Description}: The definition of the benchmarking task. \\
& \textit{Codebook question}: How is the task defined? (1-2 sentences) \\
\\
\multicolumn{2}{l}{\textbf{task\_face\_validity}} \\
& \textit{Description}: An assessment of the face validity of the benchmark. \\
& \textit{Codebook question}: Is there prima facie reason to believe that this task could benchmark the target phenomenon? \\
\\
\multicolumn{2}{l}{\textbf{task\_face\_validity\_clean}} \\
& \textit{Description}: Standardised mapping of task\_face\_validity values for statistical analysis. \\
& \textit{Codebook question}: None \\
\\
\multicolumn{2}{l}{\textbf{task\_item\_definition}} \\
& \textit{Description}: The definition and/or an example of a single item in the task. \\
& \textit{Codebook question}: What does a single item in the task dataset look like? (If the task is stored as a table, what is represented by one row in the table?) (1-2 sentences) \\
\\
\multicolumn{2}{l}{\textbf{task\_definition\_detail}} \\
& \textit{Description}: Any additional notes about the task definition. \\
& \textit{Codebook question}: Any additional details on task definition. (optional, only if something stands out) \\
\\
\multicolumn{2}{l}{\textbf{task\_source}} \\
& \textit{Description}: The source of the task items. \\
& \textit{Codebook question}: What is the source of the dataset task items? (Choose all that apply. If additional comments are needed, use the next question.) \\
\\
\multicolumn{2}{l}{\textbf{task\_source\_clean}} \\
& \textit{Description}: Standardised mapping of task\_source values for statistical analysis. \\
& \textit{Codebook question}: None \\
\\
\multicolumn{2}{l}{\textbf{task\_source\_detail}} \\
& \textit{Description}: Any additional notes about the task source. \\
& \textit{Codebook question}: Other useful notes on task source (optional, use this is something needs to be clarified) \\
\\
\multicolumn{2}{l}{\textbf{task\_ecology}} \\
& \textit{Description}: How closely does the benchmarking task resemble the real application? \\
& \textit{Codebook question}: Is the task ecologically valid? (e.g. would a person really use a model in this way?) In the case of benchmarks which cover foundational abilities across many potential applications, you may need to select multiple responses and clarify below. \\
\\
\multicolumn{2}{l}{\textbf{task\_ecology\_clean}} \\
& \textit{Description}: Standardised mapping of task\_ecology values for statistical analysis. \\
& \textit{Codebook question}: None \\
\\
\multicolumn{2}{l}{\textbf{task\_ecology\_detail}} \\
& \textit{Description}: Any additional detail about the ecological validity of the task \\
& \textit{Codebook question}: Any additional detail about the ecological validity of the task \\
\\
\multicolumn{2}{l}{\textbf{task\_train\_val}} \\
& \textit{Description}: The data splits that are provided. \\
& \textit{Codebook question}: Which of the following dataset splits are provided? (if no splits are provided, assume the entire task is the test set) \\
\\
& \textit{Summary of values:} \\
& Test: 275 \\
& Test, Train, Validation: 96 \\
& Test, Train: 51 \\
& Test, Validation: 17 \\
& Other: 2 \\
\\
\multicolumn{2}{l}{\textbf{task\_dataset\_size}} \\
& \textit{Description}: The numbers of items in the task test dataset. \\
& \textit{Codebook question}: Size of the task dataset (count, test set only, if none is reported write "NA") \\
\\
\multicolumn{2}{l}{\textbf{task\_dataset\_size\_extra}} \\
& \textit{Description}: The number of items in the task train and validation datasets, if they exist. \\
& \textit{Codebook question}: The size of the train and validation splits, if they are provided \\
\\
\multicolumn{2}{l}{\textbf{task\_dataset\_size\_detail}} \\
& \textit{Description}: Any additional notes about the task dataset size. \\
& \textit{Codebook question}: Any additional notes (e.g. the test set for some of the subcategories is very small) \\
\\
\multicolumn{2}{l}{\textbf{task\_dataset\_metadata}} \\
& \textit{Description}: Whether additional metadata is provided about the task items. \\
& \textit{Codebook question}: Does the dataset provide any metadata? (e.g. topic area, difficulty level. Do not look in the dataset, this must be described in the paper to count.) \\
\\
& \textit{Summary of values:} \\
& Yes: 322 \\
& No: 126 \\
\\
\multicolumn{2}{l}{\textbf{dataset\_metadata\_detail}} \\
& \textit{Description}: A description of any metadata provided. \\
& \textit{Codebook question}: If metadata is provided, what is it? (comma-separated list of fields, e.g. human difficulty, date, language) \\
\\
\multicolumn{2}{l}{\textbf{dataset\_sampling\_method}} \\
& \textit{Description}: The method by which task items were selected from the space of possible task items. \\
& \textit{Codebook question}: How does the dataset relate to the population it represents? (Choose all that apply, see the image for examples) \\
\\
\multicolumn{2}{l}{\textbf{dataset\_sampling\_method\_clean}} \\
& \textit{Description}: Standardised mapping of dataset\_sampling\_method values for statistical analysis. \\
& \textit{Codebook question}: None \\
\\
\multicolumn{2}{l}{\textbf{response\_format}} \\
& \textit{Description}: The format of the expected response. \\
& \textit{Codebook question}: What is the format of the expected response? (Choose all that apply. Try to stick with the provided categories, and use the next question to clarify.) \\
\\
\multicolumn{2}{l}{\textbf{response\_format\_clean}} \\
& \textit{Description}: Standardised mapping of response\_format values for statistical analysis. \\
& \textit{Codebook question}: None \\
\\
\multicolumn{2}{l}{\textbf{response\_format\_detail}} \\
& \textit{Description}: Any additional notes about the response format. \\
& \textit{Codebook question}: Any additional details about the required response format to clarify how it fits in the categories above (optional) \\
\\
\end{longtable}}

\subsection{\textcolor{metric}{Metric}}
{\small\begin{longtable}{p{.005\textwidth}p{.995\textwidth}}
\multicolumn{2}{l}{\textbf{metric\_definition}} \\
& \textit{Description}: The definition of the metric used to score the benchmark. \\
& \textit{Codebook question}: What is the primary metric for scoring the benchmark? (Choose all that apply. Please try to stick to the provided categories and elaborate below.) \\
\\
\multicolumn{2}{l}{\textbf{metric\_definition\_clean}} \\
& \textit{Description}: Standardised mapping of metric\_definition values for statistical analysis. \\
& \textit{Codebook question}: None \\
\\
\multicolumn{2}{l}{\textbf{metric\_access}} \\
& \textit{Description}: Whether the metric requires model access or not. \\
& \textit{Codebook question}: Does this metric require model access, or can it be computed from responses alone? \\
\\
& \textit{Summary of values:} \\
& Outputs alone: 422 \\
& Model access required (e.g. logits): 32 \\
\\
\multicolumn{2}{l}{\textbf{metric\_definition\_detail}} \\
& \textit{Description}: Additional details on metric definition. \\
& \textit{Codebook question}: Any additional details on metric definition. (optional, only if something stood out) \\
\\
\multicolumn{2}{l}{\textbf{metric\_face\_validity}} \\
& \textit{Description}: An assessment of the face validity of the metric. \\
& \textit{Codebook question}: Is there prima facie reason to believe that this metric could benchmark the target phenomenon? \\
\\
\multicolumn{2}{l}{\textbf{metric\_face\_validity\_clean}} \\
& \textit{Description}: Standardised mapping of metric\_face\_validity values for statistical analysis. \\
& \textit{Codebook question}: None \\
\\
\multicolumn{2}{l}{\textbf{metric\_aggregation}} \\
& \textit{Description}: The method(s) used to aggregate metric scores. \\
& \textit{Codebook question}: How are the results aggregated, if at all? (e.g. mean, weighted mean, correlation) \\
\\
\multicolumn{2}{l}{\textbf{metric\_subscores}} \\
& \textit{Description}: Whether subscores are provided for any specific subsets of the task. \\
& \textit{Codebook question}: Are scores provided for any specific subsets of the task? (e.g. by difficulty) \\
\\
& \textit{Summary of values:} \\
& Yes: 363 \\
& No: 91 \\
\\
\multicolumn{2}{l}{\textbf{metric\_subscores\_detail}} \\
& \textit{Description}: Standardised mapping of metric\_subscores values for statistical analysis. \\
& \textit{Codebook question}: If so, what subsets are provided? \\
\\
\multicolumn{2}{l}{\textbf{metric\_metascoring}} \\
& \textit{Description}: Whether the scoring involves meta-scoring techniques (pass@k, consensus@k, etc.). \\
& \textit{Codebook question}: Does the scoring involve any meta-scoring techniques? If so, which ones? \\
\\
\multicolumn{2}{l}{\textbf{metric\_fewshot}} \\
& \textit{Description}: Whether evaluation uses few-shot prompting. \\
& \textit{Codebook question}: Does the scoring involve few-shot prompting or other similar in-context learning techniques? \\
\\
& \textit{Summary of values:} \\
& No: 214 \\
& Yes: 79 \\
\\
\multicolumn{2}{l}{\textbf{metric\_statistics}} \\
& \textit{Description}: The statistics used to aggregate and compare metric scores. \\
& \textit{Codebook question}: What statistical methods are used to aggregate and compare the results? (e.g. simple mean/sum, mean and variance, clustered standard deviations) \\
\\
\multicolumn{2}{l}{\textbf{metric\_statistics\_clean}} \\
& \textit{Description}: Standardised mapping of metric\_statistics values for statistical analysis. \\
& \textit{Codebook question}: None \\
\\
\end{longtable}}

\subsection{Results and \textcolor{claims}{Claims}}
{\small\begin{longtable}{p{.005\textwidth}p{.995\textwidth}}
\multicolumn{2}{l}{\textbf{result\_interpretation}} \\
& \textit{Description}: Connection between claims and phenomenon definition. \\
& \textit{Codebook question}: Are the claims made in the results consistent with the scope of the definition being used? \\
\\
& \textit{Summary of values:} \\
& Yes: 435 \\
& No: 18 \\
\\
\multicolumn{2}{l}{\textbf{results\_comparison}} \\
& \textit{Description}: Comparison of results to other benchmarks. \\
& \textit{Codebook question}: Are comparisons made to results on other benchmarks of similar phenomena? (this requires a comparison of the nature of the results, not just a literature review) \\
\\
& \textit{Summary of values:} \\
& No: 294 \\
& Yes: 160 \\
\\
\multicolumn{2}{l}{\textbf{results\_comparison\_explanation}} \\
& \textit{Description}: Free-form explanation of the comparison. \\
& \textit{Codebook question}: If so, are theories offered to explain the similarities and differences? \\
\\
& \textit{Summary of values:} \\
& No comparisons made: 261 \\
& Yes: 144 \\
& No: 27 \\
\\
\multicolumn{2}{l}{\textbf{results\_human\_baseline}} \\
& \textit{Description}: Whether model results are compared to human performance. \\
& \textit{Codebook question}: Does the paper present a human baseline on the task? \\
\\
& \textit{Summary of values:} \\
& No: 305 \\
& Yes: 146 \\
\\
\multicolumn{2}{l}{\textbf{results\_author\_validity}} \\
& \textit{Description}: Whether benchmark authors directly address construct validity. \\
& \textit{Codebook question}: Do the authors present their own assessment of the validity of their benchmark? (i.e. do they directly address the question of construct validity for their benchmark?) \\
\\
\multicolumn{2}{l}{\textbf{results\_author\_validity\_clean}} \\
& \textit{Description}: Standardised mapping of results\_author\_validity values for statistical analysis. \\
& \textit{Codebook question}: None \\
\\
\multicolumn{2}{l}{\textbf{results\_author\_validity\_detail}} \\
& \textit{Description}: Free-form explanation of results\_author\_validity. \\
& \textit{Codebook question}: If so, please describe their evidence. \\
\\
\multicolumn{2}{l}{\textbf{results\_realism}} \\
& \textit{Description}: Whether benchmark results compare to real settings. \\
& \textit{Codebook question}: Are comparisons made between the benchmark results and results from more realistic settings? (e.g. MedQA vs supporting doctors in practice) \\
\\
& \textit{Summary of values:} \\
& No: 308 \\
& The benchmark is itself realistic: 119 \\
& Yes: 23 \\
& Other: 4 \\
\\
\multicolumn{2}{l}{\textbf{results\_realism\_clean}} \\
& \textit{Description}: Standardised mapping of results\_realism values for statistical analysis. \\
& \textit{Codebook question}: None \\
\\
\end{longtable}}

\subsection{Procedural}
{\small\begin{longtable}{p{.005\textwidth}p{.995\textwidth}}
\multicolumn{2}{l}{\textbf{authorship}} \\
& \textit{Description}: Authorship composition (industry or academia). \\
& \textit{Codebook question}: Authorship composition of the article \\
\\
& \textit{Summary of values:} \\
& Academia: 230 \\
& Mix (multiple authors from industry and academia): 186 \\
& Industry: 33 \\
& Other: 4 \\
\\
\multicolumn{2}{l}{\textbf{benchmark\_availability}} \\
& \textit{Description}: Benchmark online availability. \\
& \textit{Codebook question}: Whether the benchmark artefact is publicly available \\
\\
\multicolumn{2}{l}{\textbf{benchmark\_location}} \\
& \textit{Description}: Benchmark URL. \\
& \textit{Codebook question}: A link to the benchmark, if available (GitHub or similar) \\
\\
\multicolumn{2}{l}{\textbf{procedural\_extra}} \\
& \textit{Description}: Optional additional notes on procedural details. \\
& \textit{Codebook question}: Other useful notes on procedural details  (optional, only if something stood out) \\
\\
\multicolumn{2}{l}{\textbf{notes\_extra}} \\
& \textit{Description}: Final optional notes. \\
& \textit{Codebook question}: Any final notes about the paper not covered by above sections \\
\\
\end{longtable}}
